\subsection{可消去元素}

定义 4. 给定集合 E 上的一个合成律 \(\top\) ，映射 \(x \mapsto a \top x\) （相应地 \(x \mapsto x \top a\) ）（将 E 映入自身）称为由元素 \(a \in E\) 所作的左平移（相应地，右平移）。

在过渡到反合成律时，左平移变为右平移，反之亦然。

设 \(\pmb{\gamma}_{a},\pmb{\delta}_{a}\) （或 \(\pmb {\gamma}(a),\pmb {\delta}(a)\) ）分别表示由 \(a\in \mathbf{E}\) 所作的左平移和右平移；于是

\[
\boldsymbol {\gamma} _ {a} (x) = a \top x, \quad \boldsymbol {\delta} _ {a} (x) = x \top a.
\]

命题 1. 若律 \(\top\) 是结合的，则对所有 \(x \in E\) 和 \(y \in E\)

\[
\gamma_ {x \top y} = \gamma_ {x} \circ \gamma_ {y}, \quad \delta_ {x \top y} = \delta_ {y} \circ \delta_ {x}.
\]

对于所有 \(z \in \mathbf{E}\) :

\[
\boldsymbol {\Upsilon} _ {x \top y} (z) = (x \top y) \top z = x \top (y \top z) = \boldsymbol {\Upsilon} _ {x} (\boldsymbol {\Upsilon} _ {y} (z))
\]

\[
\boldsymbol {\delta} _ {x \top y} (z) = z \top (x \top y) = (z \top x) \top y = \boldsymbol {\delta} _ {y} (\boldsymbol {\delta} _ {x} (z))
\]

换言之，映射 \(x \mapsto \gamma_x\) 是广群 E 到映射集合 \(\mathbf{E}^{\mathbb{E}}\) （ E 到自身的映射所成的集合，带律 \((f, g) \mapsto f \circ g\) ）的一个同态；映射 \(x \mapsto \delta_x\) 是 E 到带反合成律的集合 \(\mathbf{E}^{\mathbb{E}}\) 的一个同态。若 E 是幺半群，这些同态是保单位的。

定义5. 广群 E 的元素 a 称为左（相应地，右）可消去的（或正则的），如果由 a 所作的左（相应地，右）平移是单射的。既左可消去又右可消去的元素称为可消去的（或正则的）元素。

换句话说， \(a\) 在律 \(\top\) 下可消去的充要条件是：关系式 \(a \top x = a \top y\) , \(x \top a = y \top a\) 中的每一个都蕴含 \(x = y\) （即称“a 可以从这些等式中消去”）。若存在单位元 \(e\) （在律 \(\top\) 下），则它在该律下可消去：此时平移 \(\gamma_e\) 和 \(\delta_e\) 就是 E 到自身的恒等映射。

例子。(1) 每个自然数在加法下可消去；每个自然数 \(\neq 0\) 在乘法下可消去。

\begin{enumerate}
\def\labelenumi{(\arabic{enumi})}
\setcounter{enumi}{1}
\tightlist
\item
  在一个格中，在上确界的律下除单位元（最小元）（若存在）外不存在可消去元；对下确界的律同理。特别地，在集合 E 的子集所成的集合中， \(\varnothing\) 是律 \(\cup\) 下唯一可消去的元素，而 E 是律 \(\cap\) 下唯一可消去的元素.
\end{enumerate}

命题2. 结合广群的可消去（相应地，左可消去，相应地，右可消去）元素之集构成一个子广群。

若 \(\gamma_{x}\) 和 \(\gamma_{y}\) 是单射，则 \(\gamma_{x\top y} = \gamma_{x}\circ \gamma_{y}\) 也是单射（命题1）。对 \(\delta_{x\top y}\) 类似.

